% ED26-REQ-N-c-002-02: Fagin 1976, original sections 1--4.

\subsection{Fagin 1976: finite-model probabilities and extension theory}
\label{sec:p09-n-c-002-02}

For a fixed first-order sentence, Fagin's question is whether truth settles
as the size of a uniformly sampled finite structure grows.  His answer
uses a theory of one-point extensions, first for labelled structures and
then for uniformly sampled isomorphism types.  These are different
probability spaces.  The presentation here follows the original
\emph{Probabilities on Finite Models}, sections~1--4, including the
additional remarks on its final page \cite[pp.~50--58]{Fagin1976}.
It is not the existential-second-order characterization of NP from the
separate 1974 paper, nor Libkin's later computable random-graph construction.

\begin{hbdefinition}{def:p09-fagin1976-probabilities}{The labelled and unlabelled spaces}
Let \(Y\) be a fixed finite set of relation symbols of positive arity,
with equality in the logic but with no constant or function symbols.
Write \([n]=\{1,\ldots,n\}\), for \(n\geq1\), and let \(A_n(Y)\)
be all \(Y\)-structures on \([n]\).  If \(k_i\) is the number of
\(i\)-ary symbols, then
\[
 a_n(Y)=\lvert A_n(Y)\rvert=2^{\sum_{i\geq1}k_i n^i}.
\]
Each ordered tuple, including tuples with repeated entries, independently
decides membership in each relation with probability \(1/2\).  Thus an
unrestricted binary relation has independent loops and independent
opposite directed pairs; symmetry and irreflexivity are not imposed.
Let \(\mu_n^Y\) give every member of \(A_n(Y)\) weight \(1/a_n(Y)\),
and, for a first-order \(Y\)-sentence \(\sigma\), put
\[
 p_n(\sigma)=\mu_n^Y(\{\mathcal A:\mathcal A\entails\sigma\}).
\]
This sentence probability is unchanged by adding unused relation symbols,
because every old structure has the same number of expansions.

Let \(B_n(Y)\subseteq A_n(Y)\) contain one representative of each
isomorphism type, and put \(b_n(Y)=\lvert B_n(Y)\rvert\).  The
unlabelled probability is
\[
 v_n^Y(\sigma)=
 \frac{\lvert\{\mathcal A\in B_n(Y):\mathcal A\entails\sigma\}\rvert}
      {b_n(Y)}.
\]
It is independent of the choice of representatives, since sentences are
isomorphism-invariant, but need not be independent of unused vocabulary.
An isomorphism type has \(n!/\lvert\operatorname{Aut}(\mathcal A)\rvert\)
labellings: uniform labelled sampling therefore is not uniform sampling of
types.  The superscript on \(v_n^Y\) retains this vocabulary dependence.
For example, at \(n=2\), \(v_2^{\{U\}}(\forall x\,U(x))=1/3\),
but \(v_2^{\{U,V\}}(\forall x\,U(x))=3/10\): with two unary
predicates there are ten vectors of four type counts summing to two,
and three of them put both points in the two types with \(U\) true.
\end{hbdefinition}

For finite counting with named elements, the source temporarily adds
constant symbols \(1,2,\ldots\).  In the inessential expansion of a
structure on \([n]\), the symbol \(i\) denotes \(i\) for \(i\leq n\)
and denotes \(n\) otherwise.  For example, for one binary \(P\),
\(p_n(P12)=1/2\) when \(n\geq2\) and
\(p_n(P21\land P23\land P31)=1/8\) when \(n\geq3\).
These auxiliary names are a way of writing events in a fixed probability
space, not permission to put freely interpreted constants in the main
zero--one-law vocabulary \cite[pp.~50--51]{Fagin1976}.

\begin{hbdefinition}{def:p09-fagin1976-extension}{Complete open descriptions and the theory \(T\)}
Fix distinct variables \(x_1,\ldots,x_m\), with \(m\geq0\).
A complete open description \(M(\bar x)\) selects exactly one of
\(R(z_1,\ldots,z_r)\) and its negation for every \(r\)-ary \(R\in Y\)
and every ordered tuple of variables from \(\bar x\), repetitions allowed.
There are no equality atoms in \(M\).  An extending description
\(N(\bar x,y)\) makes the corresponding choices on the enlarged variable
set and retains every conjunct of \(M\).
Define
\[
\begin{split}
 e_{M,N}\defeq\forall\bar x\biggl[
  \biggl(\bigwedge_{i<j}x_i\ne x_j\land M(\bar x)\biggr)
  \to\exists y\biggl(\bigwedge_i y\ne x_i\land N(\bar x,y)\biggr)
 \biggr].
\end{split}
\]
The theory \(T\) consists of all such \(Y\)-sentences.  Distinctness
belongs to the displayed guards, not to the open descriptions.  The case
\(m=0\) requests a point of each possible one-point description.
Empty conjunctions mean true.  \(T\) is a set of first-order sentences,
not a probability measure or an additional relation symbol.
\end{hbdefinition}

\begin{hblemma}{thm:p09-fagin1976-extension-estimate}{proved=annex:p9.n-c-002.2.extension-estimate}{The general-arity extension estimate}
For the fixed axiom \(e_{M,N}\), put
\[
 d_Y(m)=\sum_{i\geq1}k_i\bigl((m+1)^i-m^i\bigr),
 \qquad q=2^{-d_Y(m)}.
\]
For \(n>m\),
\[
 p_n(\neg e_{M,N})\leq n^m(1-q)^{n-m}\longrightarrow0.
\]
In particular, for one unrestricted binary symbol the new atomic facts
are the \(m\) facts \(P(x_i,y)\), the \(m\) facts \(P(y,x_i)\),
and \(P(y,y)\), giving \(q=2^{-(2m+1)}\).
The general formula counts every tuple involving the new point, including
several occurrences of it; it is not a reduction to a binary signature.
\end{hblemma}

\begin{hbtheorem}{thm:p09-fagin1976-theory}{proved=annex:p9.n-c-002.2.theory}{Existence, back-and-forth, and completeness}
The theory \(T\) is consistent, has no finite models, and has exactly
one countably infinite model up to isomorphism.  It is complete: for every
first-order \(Y\)-sentence \(\sigma\), exactly one of
\(T\entails\sigma\) and \(T\entails\neg\sigma\) holds.
\end{hbtheorem}

Fagin attributes the theory and this theorem to Gaifman.  The original
gives a construction suggestion and a back-and-forth sketch, invoking the
\L o\'s--Vaught test.  The annex expands the finite partial-isomorphism
extension, both directions of the enumeration, and the passage from
countable categoricity to completeness.  It uses ordinary first-order
compactness and downward L\"owenheim--Skolem as mathematical prerequisites;
it does not claim that Fagin printed these expanded proofs
\cite[thm.~1, p.~51]{Fagin1976}.

\begin{hbtheorem}{thm:p09-fagin1976-labelled}{proved=annex:p9.n-c-002.2.labelled}{The labelled zero--one law and its rate}
Every first-order \(Y\)-sentence \(\sigma\) has a limiting labelled
probability in \(\{0,1\}\), and
\[
 T\entails\sigma\quad\Longleftrightarrow\quad p_n(\sigma)\longrightarrow1.
\]
The convergence is geometrically fast in the source's sense: if the
limit is \(\ell\), there are \(0<\theta<1\) and \(N\) such that
\(\lvert p_n(\sigma)-\ell\rvert<\theta^n\) for every \(n>N\).
The constants depend on the fixed sentence and vocabulary.
\end{hbtheorem}

The logical bridge matters: probability one for each individual extension
axiom does not itself give probability one for their infinite conjunction.
Compactness supplies a \emph{finite} set of axioms entailing the fixed
sentence, and a finite union bound supplies its asymptotic probability.
Completeness then supplies the dichotomy.  This reconstructs Theorem~2,
Corollary~3, Theorem~4, and the labelled part of Theorem~10
\cite[pp.~52--53 and 56]{Fagin1976}.

\begin{hbcorollary}{thm:p09-fagin1976-finite-fragments}{proved=annex:p9.n-c-002.2.finite-fragments}{Finite fragments and the spectrum qualification}
Every finite \(T_0\subseteq T\) has a model of cardinality \(n\) for
every sufficiently large \(n\).  For any first-order sentence \(\sigma\)
in a finite vocabulary, even allowing positive-arity functions and
constants, either \(\sigma\) or \(\neg\sigma\) has a cofinite spectrum,
where the spectrum is the set of positive finite model cardinalities.
This does not assert that the spectrum of each sentence is finite or
cofinite, or that functions and constants satisfy the probability law.
\end{hbcorollary}

The source gives two arguments for the spectrum assertion.  Its first
mentions replacing functions by relations to preserve a sentence's
spectrum.  Such a replacement alone does not justify transferring the
spectrum of the \emph{negated} replacement back to \(\neg\sigma\):
the negated replacement can also hold because the coding relations fail
their function constraints.  The annex therefore uses the source's second
argument, in which all relations are full, all functions project onto their
first coordinate, and all constants coincide in a canonical infinite
structure.  This supplies the assertion for the original sentence without
an unjustified negation transfer \cite[cor.~6 and its two proofs, p.~53]{Fagin1976}.

For a relational \(\sigma\) with \(p_n(\sigma)\to1\), each fixed
\(0<\beta<1\) eventually gives at least
\(\beta a_n(Y)/n!\) distinct isomorphism types of models of size \(n\).
There are at least \(\beta a_n(Y)\) labelled models and at most \(n!\)
labellings per type.  If a relation has arity at least two, then
\(a_n(Y)\geq2^{n^2}\) and
\(\log_2 n!\leq n\log_2 n\); thus this lower bound grows without bound.
This is the ``many nonisomorphic models'' observation at the end of
section~2, not yet the equality of the two probability limits.

\subsubsection{Uniform isomorphism types: two different cases}

\begin{hblemma}{thm:p09-fagin1976-monadic}{proved=annex:p9.n-c-002.2.monadic}{Monadic counts and first-order thresholds}
Suppose \(Y\) consists of \(h\) unary predicates and put \(s=2^h\).
Its one-point types are the \(s\) choices of which predicates hold.
An isomorphism type on \([n]\) is uniquely described by nonnegative
counts \((c_1,\ldots,c_s)\) with sum \(n\).  Consequently
\[
 b_n(Y)=\binom{n+s-1}{s-1}.
\]
When \(s\geq2\), the unlabelled probability that a specified point
type has exactly \(k\) members, for \(n>k\), is
\[
 \frac{\binom{n-k+s-2}{s-2}}{\binom{n+s-1}{s-1}}
 \longrightarrow0.
\]
A sentence of quantifier rank \(r\) cannot distinguish two monadic
structures whose point-type counts agree after truncation at \(r\).
Here quantifier rank is the maximum nesting depth of quantifiers:
atoms have rank zero, Boolean combinations take the maximum of their
operand ranks, and each quantifier increases rank by one.
It follows that every monadic first-order sentence is a Boolean
combination of assertions of exactly \(k\) members of a point type,
with \(k<r\).  Its unlabelled limit is zero or one and agrees with
its labelled limit.  In fact \(v_n^Y(\sigma)\) is eventually a ratio
of polynomials in \(n\).  For \(h=0\), there is just one type and
one structure per cardinality; equality-only sentences are eventually
constant, so the same limit conclusion holds.
\end{hblemma}

This supplies the stars-and-bars lemma and the monadic normal-form step
used in the source's Lemma~7 and first case of Theorem~9, rather than
silently applying a nonmonadic enumeration estimate to unary predicates
\cite[pp.~54--55]{Fagin1976}.

\begin{hbtheorem}{thm:p09-fagin1976-enumeration}{cited=Fagin1976}{The imported nonmonadic enumeration estimate}
If \(Y\) contains a relation of arity at least two, then
\[
 \frac{b_n(Y)n!}{a_n(Y)}\longrightarrow1.
\]
Moreover \(a_n(Y)/(b_n(Y)n!)\) converges to one geometrically fast.
The first assertion is Fagin's Theorem~8; the rate is invoked in the
proof of Theorem~10.  Both are imported there from reference~[3],
\emph{The number of finite relational structures}, IBM research report
RC5587, August 1975.  This entry uses that cited estimate and does not
claim to reconstruct the separate report's enumeration proof.
\end{hbtheorem}

\begin{hbtheorem}{thm:p09-fagin1976-unlabelled}{proved=annex:p9.n-c-002.2.unlabelled}{Unlabelled transfer, with its enumeration dependency}
For every first-order \(Y\)-sentence,
\[
 \lim_{n\to\infty}v_n^Y(\sigma)=\lim_{n\to\infty}p_n(\sigma)\in\{0,1\}.
\]
In the nonmonadic case the convergence is geometrically fast, using
\cref{thm:p09-fagin1976-enumeration}.  This rate need not hold in the
monadic case.  For one unary predicate \(U\),
\[
 v_n^{\{U\}}(\forall x\,U(x))=\frac1{n+1},
 \qquad p_n(\forall x\,U(x))=2^{-n}.
\]
For one unrestricted binary relation \(P\), in contrast,
\(p_n(\forall x\,P(x,x))=2^{-n}\), and its unlabelled probability
also converges geometrically to zero.
\end{hbtheorem}

\subsubsection{The graph link and the limits of changing the measure}

\begin{hbtheorem}{thm:p09-fagin1976-graph}{proved=annex:p9.n-c-002.2.graph}{A fixed-density graph extension calculation}
In this result only, replace unrestricted relations by simple graphs:
one binary edge symbol \(E\), with \(E(x,x)\) false and
\(E(x,y)\leftrightarrow E(y,x)\).  Independently for each unordered
pair, choose an edge with fixed probability \(p\), where \(0<p<1\)
does not depend on \(n\).  For disjoint fixed vertex sets \(U,V\)
of sizes \(u,v\), a fresh vertex is required to meet all of \(U\)
and none of \(V\).  Its success probability is
\(q=p^u(1-p)^v>0\), and
\[
 \Pr[\hbox{some instance with these sizes has no fresh witness}]
 \leq n^{u+v}(1-q)^{n-u-v}\longrightarrow0.
\]
The theory comprising symmetry, irreflexivity, and every such extension
sentence is consistent and complete by back-and-forth.  Every fixed
first-order graph sentence has limiting probability zero or one,
independently of which fixed \(p\in(0,1)\) is used.  Limit one means
truth in the unique countably infinite model of this extension theory.
\end{hbtheorem}

Explicitly, the extension sentence for \(m=u+v\) is
\[
\begin{split}
 \forall x_1\cdots x_m\biggl[\bigwedge_{i<j}x_i\ne x_j
 \to\exists y\biggl(&\bigwedge_{i=1}^m y\ne x_i
     \land\bigwedge_{i=1}^u E(y,x_i)\\
     &{}\land\bigwedge_{i=u+1}^m\neg E(y,x_i)\biggr)\biggr].
\end{split}
\]
The probability estimate concerns \(n>m\); the axiom itself also
applies to smaller nonempty graph structures with its displayed guard.

For example, with \(u=v=1\) and \(p=1/2\), each candidate succeeds
with probability \(1/4\), and the failure bound is
\(n^2(3/4)^{n-2}\).  With \(u=1,v=0\), the extension axioms give
\(\forall x\exists y\,E(x,y)\) limiting probability one; hence
\(\exists x\forall y\neg E(x,y)\), the assertion of an isolated
vertex, has limit zero.  The countable model is a model of the limiting
theory, not a set-theoretic limit obtained by joining the finite graphs.

For a sentence \(\tau\) with \(p_n(\tau)>0\), define
\(p_n(\sigma\mid\tau)=p_n(\sigma\land\tau)/p_n(\tau)\).
Define \(v_n^Y(\sigma\mid\tau)\) by the same ratio for that
measure, when its denominator is positive.  Fagin's section~4(6)
obtains the \(p=1/2\) graph model by conditioning
the unrestricted binary model on
\[
 \tau_1=\forall x\neg P(x,x)\land
             \forall x\forall y(P(x,y)\leftrightarrow P(y,x)).
\]
Each simple labelled graph corresponds to exactly one conditioned
structure, so unordered edges remain independent fair choices.  The
source expresses this correspondence using a graph predicate \(Q\):
replace each occurrence of \(P\) in \(\sigma\) by \(Q\) to
obtain \(\sigma^Q\).  Then \(p_n(\sigma\mid\tau_1)\) is
exactly the probability of \(\sigma^Q\) in the fair graph space.
The displayed \(E\)-presentation is the same graph language with
its edge symbol renamed, not an extra relation.  The
annex also supplies his tournament variant, conditioning on
\[
 \tau_2=\forall x\neg P(x,x)\land
  \forall x\forall y\bigl(x\ne y\to
             (P(x,y)\leftrightarrow\neg P(y,x))\bigr),
\]
and redefining complete descriptions to respect one oriented edge per
unordered pair \cite[p.~57]{Fagin1976}.  These constrained models are
not the original uniform space of all binary relations.

\begin{hbtheorem}{thm:p09-fagin1976-boundaries}{proved=annex:p9.n-c-002.2.boundaries}{Source counterexamples and conditional probability}
The following limitations refer to the explicitly specified measures.
The abbreviation \(\exists!y\,\phi(y)\) means that exactly one
point satisfies \(\phi\):
\(\exists y(\phi(y)\land\forall z(\phi(z)\to z=y))\), with
\(z\) fresh and substitution capture-free.
Over the empty vocabulary, the second-order sentence
\[
 \exists P\,\forall x\exists!y\,
        (x\ne y\land P(x,y)\land P(y,x))
\]
holds exactly at even positive cardinalities.  Both its labelled and
unlabelled truth probabilities therefore oscillate between zero and one.
For uniform labelled structures with one unary predicate \(U\) and
one freely interpreted constant \(c\), \(p_n(U(c))=1/2\).
For one unary function \(F\),
\[
 p_n(\forall x\,F(x)\ne x)=\frac{(n-1)^n}{n^n}
                         \longrightarrow e^{-1}.
\]
These are counterexamples to the zero--one conclusion, not proofs that
every function-vocabulary sentence lacks a limit.

Even conditioning on a first-order sentence that has positive probability
at every cardinality need not yield a limit.  Put
\[
\begin{split}
 \sigma_1&=\forall x\exists!y\,(x\ne y\land P(x,y)\land P(y,x)),\\
 \sigma_2&=\exists w\,\forall x\bigl(x\ne w\to
    \exists!y\,(y\ne w\land x\ne y\land P(x,y)\land P(y,x))\bigr),\\
 \tau&=\sigma_1\lor\sigma_2.
\end{split}
\]
Then \(p_n(\tau)>0\) and \(v_n^{\{P\}}(\tau)>0\) for every
\(n\geq1\), but both conditional probabilities of \(\sigma_1\)
given \(\tau\) are one at even \(n\) and zero at odd \(n\).
By contrast, for \(\tau_3=\forall x\forall y(U(x)\leftrightarrow U(y))\),
each first-order conditional probability \(p_n(\sigma\mid\tau_3)\)
is eventually one of \(0,1/2,1\); the universal-\(U\) sentence has
conditional probability \(1/2\), not a zero--one limit.
\end{hbtheorem}

The source asks open questions about convergence with functions and
constants and about which conditionings ensure convergence,
including whether a universal conditioning sentence always ensures
convergence \cite[pp.~56--58]{Fagin1976}.  These are questions recorded in the
1976 article, not assertions of their present research status.

\begin{hbwarning}{def:p09-fagin1976-outside-hypotheses}{Order, nullary predicates, and changing edge density}
A nullary random predicate has sentence probability \(1/2\); the
positive-arity convention is material.  A fixed built-in linear order
with independent fair unary \(U\) gives the first-order sentence
``the least element has \(U\)'' probability \(1/2\).  This is an
ordered, constrained space, not uniform sampling of all interpretations
of an unconstrained order symbol.  If simple-graph edge probability
instead varies as \(p(n)=1/n^2\), the first-order sentence asserting no
edge has probability, for \(n\geq2\),
\((1-1/n^2)^{n(n-1)/2}\to e^{-1/2}\).
Thus neither arbitrary order enrichments nor varying probabilities are
licensed by the fixed-density extension calculation.  These elementary
boundary examples explain the hypotheses; they are not new claims about
all ordered structures or all sparse random-graph regimes.
\end{hbwarning}

\begin{hbtheorem}{thm:p09-fagin1976-mixed-measure}{cited=Fagin1976}{The final-page nonuniform measure, stated without proof in the source}
For one unrestricted binary relation, fix \(0<p<1\) and
\(0\leq\alpha<1\), and write \(f(t)=(1-\alpha)p^t+\alpha\).
A structure with \(r\) true and \(s\) false ordered-pair facts,
\(r+s=n^2\), receives weight
\[
 \begin{cases}
  f(n^2),&s=0,\\
  (1-f(1))p^r(1-p)^{s-1},&s>0.
 \end{cases}
\]
The weights sum to one: they are the mixture of mass \(\alpha\) on
the full relation and mass \(1-\alpha\) on independent Bernoulli
\(p\) choices for all ordered pairs.  Under this measure every
first-order sentence has a limit belonging to
\(\{0,\alpha,1-\alpha,1\}\).  In particular the probability of
\(\forall x\forall y\,P(x,y)\) tends to \(\alpha\).
This is section~4(7)'s separately specified measure and convergence
assertion; the article expressly states the convergence assertion without
proof \cite[p.~58]{Fagin1976}.
\end{hbtheorem}

The weights are transparent as a mixture: mass \(\alpha\) is placed
on the full relation and the remaining mass is independent Bernoulli
\(p\) on all ordered pairs, since
\(1-f(1)=(1-\alpha)(1-p)\).  In particular the displayed universal
sentence has probability \(\alpha+(1-\alpha)p^{n^2}\).
The source also remarks that its vocabulary can be generalized; it does
not specify all general-vocabulary weights or print the proof, so that
remark is not expanded into an unrestricted measure theorem here.

\begin{proofboundary}{def:p09-fagin1976-source-boundary}{What is proved here and what remains imported}
The annex supplies the general-arity estimate, the completed
back-and-forth/completeness argument, the labelled law, spectrum argument,
monadic normal form and counting, unlabelled transfer, graph/tournament
calculation, and the stated counterexamples.  Its unlabelled nonmonadic
proof depends on the expressly cited enumeration theorem and rate from
Fagin's reference~[3].  No proof of that separate 1975 report is claimed.
The final-page mixed-measure convergence statement retains its source's
statement-only standing.  These are the original presentation's imported
and unproved boundaries, not substitutions of another version's proof.
\end{proofboundary}
